% openany：章可起于任意页，章与章之间不插空白页（默认 openright 会在偶数页留白）
\documentclass[12pt,a4paper,openany]{book}
\usepackage{framed}
\usepackage[left=2.8cm,right=2.8cm,
top=2.8cm,bottom=2.8cm]{geometry}

\raggedbottom

\linespread{1.20}

% 段落：中文专著体例——段首缩进，段间不靠大间距区分
% （具体数值与章/节首段强制缩进见文末「段落缩进加固」）
\setlength{\parindent}{2em}
\setlength{\parskip}{0.2em}

\newenvironment{historicalnote}
  {\par\vspace{4pt}\noindent\textbf{历史说明：}\begin{framed}\small}
  {\end{framed}\par\vspace{4pt}}
%======================================================================
% 1. 基础包导入
%======================================================================
% 编码与中文字体支持
\usepackage{xeCJK}
\setCJKmainfont[ItalicFont=Kaiti SC]{Songti SC}
\setCJKmonofont{Songti SC}
% 章/节后首段缩进（英文 book 默认首段顶格；中文需 indentfirst）
\usepackage{indentfirst}

% 数学公式与符号
\usepackage{amsmath, amssymb, amsthm}
\allowdisplaybreaks
\usepackage{mathtools}
\usepackage{bm}

% 页面字体与基本设置
\usepackage{lmodern}
\usepackage{xcolor}
\usepackage{comment}
\usepackage{enumitem}

% 表格与图表标题相关包
\usepackage{array}
\usepackage{caption}
\usepackage{subcaption}
\usepackage{booktabs}
\usepackage{longtable}
% 全书图题 / 表题：Figure → 图，Table → 表
\renewcommand{\figurename}{图}
\renewcommand{\tablename}{表}
\captionsetup[figure]{name=图}
\captionsetup[table]{name=表}

% 算法伪代码相关包
\usepackage{algorithm}
\usepackage{algorithmic}

% 超链接设置
\usepackage{hyperref}
\hypersetup{
  colorlinks=true,
  linkcolor=blue!60!black,
  urlcolor=blue!60!black,
  pdftitle={素数分布与黎曼猜想},
  pdfauthor={作者},
  pdfsubject={素数分布与黎曼猜想}
}
% \autoref 显示为「图 / 表」而非 Figure / Table
\renewcommand{\figureautorefname}{图}
\renewcommand{\tableautorefname}{表}

% 附录支持（toc=在目录中显示附录条目；page=在第一个附录前插入"附录"分隔页）
\usepackage[toc,page]{appendix}
\renewcommand{\appendixtocname}{附录目录}
\renewcommand{\appendixpagename}{附\quad 录}
\renewcommand{\appendixname}{附录}
\renewcommand{\contentsname}{目录}

% 样式框包
\usepackage{mdframed}
\usepackage{tcolorbox}
\tcbuselibrary{skins,breakable,theorems}

% Pgfplots 绘图包及设置
\usepackage{pgfplots}
\pgfplotsset{compat=1.17}
\usepgfplotslibrary{groupplots}

% TikZ 图形库及常用库
\usepackage{tikz}
\usetikzlibrary{
  positioning,
  arrows.meta,
  decorations.markings,
  decorations.pathmorphing,
  calc,
  angles,
  quotes,
  shapes.multipart,
  backgrounds,
  fit
}

%======================================================================
% 2. 颜色系统定义
%======================================================================
% 基础色彩
\definecolor{MainBlue}{RGB}{0,60,120}
\definecolor{LightBlue}{RGB}{200,220,255}
\definecolor{DarkRed}{RGB}{140,20,20}
\definecolor{DarkGreen}{RGB}{10,105,45}
\definecolor{Gold}{RGB}{148,108,0}
\definecolor{Purple}{RGB}{90,0,130}
\definecolor{Gray}{RGB}{70,70,70}
\definecolor{LightRed}{RGB}{255,238,238}
\definecolor{LightGreen}{RGB}{232,248,236}
\definecolor{LightGold}{RGB}{255,250,230}
\definecolor{LightPurple}{RGB}{245,235,255}

% 方案 II：轻量色标（结构 / 辅助 / 功能 至多三类）
% 结构=冷蓝灰；辅助=中性灰；功能框=深灰细线（解读、章末汇总）
\definecolor{StructColor}{RGB}{55,90,125}    % 定义/定理/引理/命题/推论/算法
\definecolor{AuxColor}{RGB}{100,100,100}     % 例、注
\definecolor{FuncColor}{RGB}{75,75,80}       % 原稿解读等功能框
% 兼容旧色名（映射到结构/辅助，避免章内残留色名失效）
\colorlet{DefColor}{StructColor}
\colorlet{ThmColor}{StructColor}
\colorlet{LemColor}{StructColor}
\colorlet{PropColor}{StructColor}
\colorlet{CorColor}{StructColor}
\colorlet{RemColor}{AuxColor}
\definecolor{TextColor}{RGB}{178,34,34}
\definecolor{colorPrime}{RGB}{180,0,0}

%======================================================================
% 3. 页面布局与页眉页脚设置
%======================================================================
\usepackage{geometry}
\geometry{margin=1in}

\usepackage{fancyhdr}
\pagestyle{fancy}
\fancyhf{}
\fancyfoot[C]{素数分布与黎曼猜想 \quad \thepage}
\renewcommand{\headrulewidth}{0pt}
\setcounter{secnumdepth}{2}
\setcounter{tocdepth}{2}

\fancypagestyle{firstpage}{
  \fancyhf{}
  \fancyhead[L]{\textcolor{TextColor}{\Large 在读}}
  \fancyfoot[C]{复变函数基础 \quad \thepage}
  \renewcommand{\headrulewidth}{0pt}
}

%======================================================================
% 4. 数学命令与自定义算子
%======================================================================
% 常用简写与符号
\newcommand{\dif}{\,\mathrm{d}}
\newcommand{\Res}{\operatorname{Res}}
\newcommand{\re}[1]{\operatorname{Re}\!\left(#1\right)}
\newcommand{\im}[1]{\operatorname{Im}\!\left(#1\right)}
\newcommand{\Arg}{\operatorname{Arg}}
\newcommand{\Log}{\operatorname{Log}}
\newcommand{\CC}{\mathbb{C}}
\newcommand{\RR}{\mathbb{R}}
\newcommand{\ZZ}{\mathbb{Z}}

\DeclareMathOperator{\Real}{Re}
\DeclareMathOperator{\Imag}{Im}

% 数集统一定义
\DeclareMathOperator{\R}{\mathbb{R}}
\DeclareMathOperator{\N}{\mathbb{N}}
\DeclareMathOperator{\Z}{\mathbb{Z}}
\DeclareMathOperator{\Q}{\mathbb{Q}}
\DeclareMathOperator{\C}{\mathbb{C}}

% 配对括号与数学运算符
\DeclarePairedDelimiter\abs{\lvert}{\rvert}
\DeclarePairedDelimiter\floor{\lfloor}{\rfloor}

% 黎曼 Zeta 函数相关特殊符号
\newcommand{\sumrho}{\sum_{\rho}}
\DeclareMathOperator{\Li}{Li}
\DeclareMathOperator{\li}{li}
\DeclareMathOperator{\Ei}{Ei}

%======================================================================
% 5. 定理环境（方案 II：轻量色标）
%======================================================================
% 结构框 / 辅助框：无填充底（不用 colback 白底，也不用 opacityback=0——后者 Xe 会变黑）
% empty 皮肤 + borderline 只描边；标题为框内加粗黑字。
% 证明：无框；结构框可跨页。
\tcbset{
  thmbase/.style={
    empty, breakable,
    sharp corners,
    before skip=8pt, after skip=6pt,
    boxsep=0pt,
    left=6pt, right=6pt, top=4pt, bottom=4pt,
    borderline={0.45pt}{0pt}{StructColor!40!black},
    coltitle=black,
    fonttitle=\bfseries\small,
    colframe=StructColor!40!black,
    colback=white,
    separator sign={.\ },
    attach title to upper,
    title after break={（续）},
    pad before break=1mm,
    pad after break=1mm,
  },
  notebase/.style={
    empty,
    sharp corners,
    before skip=6pt, after skip=5pt,
    boxsep=0pt,
    left=6pt, right=6pt, top=3pt, bottom=3pt,
    borderline={0.35pt}{0pt}{AuxColor!30!black},
    coltitle=black,
    fonttitle=\bfseries\small,
    colframe=AuxColor!30!black,
    colback=white,
    separator sign={.\ },
    attach title to upper,
  }
}

\newtcbtheorem[number within=section]{theorem}{定理}{thmbase}{thm}
\newtcbtheorem[use counter from=theorem]{corollary}{推论}{thmbase}{cor}
\newtcbtheorem[use counter from=theorem]{lemma}{引理}{thmbase}{lem}
\newtcbtheorem[use counter from=theorem]{proposition}{命题}{thmbase}{prop}
\newtcbtheorem[use counter from=theorem]{definition}{定义}{thmbase}{def}
\newtcbtheorem[use counter from=theorem]{alg}{算法}{thmbase}{alg}
\newtcbtheorem[use counter from=theorem]{remark}{注}{notebase}{rem}
\newtcbtheorem[use counter from=theorem]{example}{例}{notebase}{ex}

% 证明环境（保留 amsthm 原生 proof，重定义为中文格式）
\renewcommand{\proofname}{证明}
\makeatletter
\renewenvironment{proof}[1][\proofname]{\par
  \pushQED{\qed}%
  \normalfont \topsep6\p@\@plus6\p@
  \trivlist
  \item[\hskip\labelsep\bfseries#1\@addpunct{：}]\ignorespaces
}{%
  \popQED\endtrivlist\@endpefalse
}
\makeatother

%======================================================================
% 6. 自定义样式框与卡片环境
%======================================================================
% 功能框：原稿「解读」—— empty + 描边，无填充
\definecolor{CommentaryFrame}{RGB}{90, 90, 95}
\definecolor{CommentaryTitle}{RGB}{50, 50, 55}
% 跨页统一：首页【解读】标题；续页【解读】接上框（勿再手拆成两个 commentary）
% 可选参数：如 \begin{commentary}[unbreakable]{标题} 强制整框不跨页
\newtcolorbox{commentary}[2][]{
  empty,
  breakable,
  sharp corners,
  boxsep=0pt,
  left=8pt, right=8pt, top=4pt, bottom=5pt,
  borderline={0.5pt}{0pt}{CommentaryFrame},
  coltitle=CommentaryTitle,
  fonttitle=\bfseries\small,
  colframe=CommentaryFrame,
  colback=white,
  before skip=8pt,
  after skip=6pt,
  title={【解读】#2},
  attach title to upper,
  title after break={【解读】接上框},
  pad before break=3mm,
  pad after break=3mm,
  #1
}

% 要点/结果框：与定理同为 empty+描边（mdframed 常自带底色，已弃用）
% 可选参数：如 \begin{resultframe}[nobreak]
\newtcolorbox{keyframe}[1][]{
  empty, breakable, sharp corners,
  boxsep=0pt,
  left=8pt, right=8pt, top=6pt, bottom=6pt,
  borderline={0.5pt}{0pt}{FuncColor!70!black},
  before skip=8pt, after skip=8pt,
  #1
}
\newtcolorbox{resultframe}[1][]{
  empty, breakable, sharp corners,
  boxsep=0pt,
  left=8pt, right=8pt, top=6pt, bottom=6pt,
  borderline={0.45pt}{0pt}{AuxColor!45!black},
  before skip=8pt, after skip=8pt,
  #1
}
\newtcolorbox{explainframe}[1][]{
  empty, breakable, sharp corners,
  boxsep=0pt,
  left=8pt, right=8pt, top=5pt, bottom=5pt,
  borderline={0.4pt}{0pt}{AuxColor!40!black},
  before skip=6pt, after skip=6pt,
  #1
}

%======================================================================
% 7. TikZ 图形库与多阶段编译设置
%======================================================================
% 编译阶段定义
\def\stage{1} % 默认编译第三阶段关系图（最终大统一版），所有概念和连线全部着色

% ---------- 阶段自适应颜色编译定义 ----------
\def\stageOne{1}
\def\stageTwo{2}
\ifx\stage\stageOne
  % === 阶段 1：仅初等数论、Gamma与Zeta上半部彩显 ===
  \colorlet{colorPrime}{blue!12}
  \colorlet{colorMultiplicative}{teal!15}
  \colorlet{colorApproximation}{yellow!25}
  \colorlet{colorComplex}{gray!10}
  \colorlet{colorRH}{gray!5}
  \colorlet{colorGammaUpper}{blue!15}
  \colorlet{colorGammaLower}{red!15}
  \colorlet{drawGamma}{black}
  \colorlet{textGammaUpper}{black}
  \colorlet{textGammaLower}{black!70}
  \colorlet{colorZetaUpper}{blue!15}
  \colorlet{colorZetaLower}{gray!10}
  \colorlet{drawZeta}{black}
  \colorlet{textZetaUpper}{black}
  \colorlet{textZetaLower}{gray!50}
  \colorlet{drawPrime}{black}
  \colorlet{drawMultiplicative}{teal!60!black}
  \colorlet{drawApproximation}{orange!70!black}
  \colorlet{drawComplex}{gray!40}
  \colorlet{drawRH}{gray!30}
  \colorlet{textApproximation}{black}
  \colorlet{textComplex}{gray!60}
  \colorlet{textRH}{gray!60}
  \colorlet{colorRHMu}{gray!60}
  \colorlet{textRHHighlight}{gray!60}
  \colorlet{lineArithmetic}{blue!70!black}
  \colorlet{lineElementary}{violet!80!black}
  \colorlet{lineDefinition}{black!80}
  \colorlet{lineComplexDef}{gray!50}
  \colorlet{lineBasicAnalysis}{green!50!black}
  \colorlet{lineAnalysis}{gray!50}
  \colorlet{lineZeros}{gray!50}
  \colorlet{lineGammaExtension}{red!70!black}
  \colorlet{lineExtension}{gray!50}
  \colorlet{linePNT}{brown!70!black}
\else
  \ifx\stage\stageTwo
    % === 阶段 2：初等数论 + 解析延拓彩色，显式公式与零点结构灰度 ===
    \colorlet{colorPrime}{blue!12}
    \colorlet{colorMultiplicative}{teal!15}
    \colorlet{colorApproximation}{yellow!25}
    \colorlet{colorComplex}{cyan!15}
    \colorlet{colorRH}{gray!5}
    \colorlet{colorGammaUpper}{blue!15}
    \colorlet{colorGammaLower}{red!15}
    \colorlet{drawGamma}{black}
    \colorlet{textGammaUpper}{black}
    \colorlet{textGammaLower}{black!70}
    \colorlet{colorZetaUpper}{blue!15}
    \colorlet{colorZetaLower}{red!15}
    \colorlet{drawZeta}{black}
    \colorlet{textZetaUpper}{black}
    \colorlet{textZetaLower}{black!70}
    \colorlet{drawPrime}{black}
    \colorlet{drawMultiplicative}{teal!60!black}
    \colorlet{drawApproximation}{orange!70!black}
    \colorlet{drawComplex}{cyan!60!black}
    \colorlet{drawRH}{gray!30}
    \colorlet{textApproximation}{black}
    \colorlet{textComplex}{black}
    \colorlet{textRH}{gray!60}
    \colorlet{colorRHMu}{gray!60}
    \colorlet{textRHHighlight}{gray!60}
    \colorlet{lineArithmetic}{blue!70!black}
    \colorlet{lineElementary}{violet!80!black}
    \colorlet{lineDefinition}{black!80}
    \colorlet{lineComplexDef}{black!80}
    \colorlet{lineBasicAnalysis}{green!50!black}
    \colorlet{lineAnalysis}{gray!50}
    \colorlet{lineZeros}{gray!50}
    \colorlet{lineGammaExtension}{red!70!black}
    \colorlet{lineExtension}{red!70!black}
    \colorlet{linePNT}{brown!70!black}
  \else
    % === 阶段 3：全谱系大统一，全部彩显 ===
    \colorlet{colorPrime}{blue!12}
    \colorlet{colorMultiplicative}{teal!15}
    \colorlet{colorApproximation}{yellow!25}
    \colorlet{colorComplex}{cyan!15}
    \colorlet{colorRH}{orange!15}
    \colorlet{colorGammaUpper}{blue!15}
    \colorlet{colorGammaLower}{red!15}
    \colorlet{drawGamma}{black}
    \colorlet{textGammaUpper}{black}
    \colorlet{textGammaLower}{black!70}
    \colorlet{colorZetaUpper}{blue!15}
    \colorlet{colorZetaLower}{red!15}
    \colorlet{drawZeta}{black}
    \colorlet{textZetaUpper}{black}
    \colorlet{textZetaLower}{black!70}
    \colorlet{drawPrime}{black}
    \colorlet{drawMultiplicative}{teal!60!black}
    \colorlet{drawApproximation}{orange!70!black}
    \colorlet{drawComplex}{cyan!60!black}
    \colorlet{drawRH}{orange!70!black}
    \colorlet{textApproximation}{black}
    \colorlet{textComplex}{black}
    \colorlet{textRH}{black}
    \colorlet{colorRHMu}{teal!60!black}
    \colorlet{textRHHighlight}{orange!70!black}
    \colorlet{lineArithmetic}{blue!70!black}
    \colorlet{lineElementary}{violet!80!black}
    \colorlet{lineDefinition}{black!80}
    \colorlet{lineComplexDef}{black!80}
    \colorlet{lineBasicAnalysis}{green!50!black}
    \colorlet{lineAnalysis}{green!50!black}
    \colorlet{lineZeros}{orange!80!black}
    \colorlet{lineGammaExtension}{red!70!black}
    \colorlet{lineExtension}{red!70!black}
    \colorlet{linePNT}{brown!70!black}
  \fi
\fi

% TikZ 箭头样式
\tikzset{marrow/.style={
  decoration={markings,mark=at position #1 with {\arrow[scale=1.3]{Stealth}}},
  postaction={decorate}},marrow/.default=0.55}

%======================================================================
% 段落缩进加固（须放在 preamble 末尾）
%======================================================================
% 1) 中文成书：章/节后首段也缩进（英文 book 默认首段顶格）
\makeatletter
\let\@afterindentfalse\@afterindenttrue
% 2) caption 包：题注在分组内把 \parindent 置 0；\caption@end 仅为 \endgroup，
%    组外不会自动恢复。本书多处图后 \captionof 会因此污染后记等后续正文。
%    必须在 \caption@end（分组结束）之后恢复。
\newcommand{\grok@restore@parindent}{%
  \setlength{\parindent}{2em}%
  \setlength{\parskip}{0.2em}%
}
\let\grok@orig@caption@end\caption@end
\renewcommand*\caption@end{%
  \grok@orig@caption@end
  \grok@restore@parindent
}
\makeatother
\setlength{\parindent}{2em}
\setlength{\parskip}{0.2em}